\section{Угрозы валидности}
\label{sec:threats}

\paragraph{Одно семейство моделей.} Каждый запуск использует
\texttt{\EOneModel}. Парный пилот на более сильной модели того же
провайдера не решил ни одной задачи сверх того, что решала маленькая
модель, — при \num{22.7}$\times$ медианной стоимости на ячейку\live, и
именно поэтому серию не стали переносить на неё. Этот пилот — всего 7--8
задач, он закрывает вопрос бюджета, а не вопрос переносимости. Мы ничего не
утверждаем о моделях других семейств и, в частности, о том, поддерживала бы
модель с более сильным следованием инструкциям состояние настолько хорошо,
чтобы избежать режима отказа из \S\ref{sec:mechanism}.

\paragraph{Один корпус и его диапазон горизонта.} Задачи Terminal-Bench 2.0
при контрольной группе используют медианно лишь около \SI{5}{\percent} окна
контекста\live. Окно поэтому редко становится связывающим ограничением — а
это как раз тот режим, где механизм памяти должен быть наименее нужен. На
более длинном Terminal-Bench 4.0 пилот показывает медианный пик уже в
\SI{66}{\percent} окна\live, и там тот же механизм проявляется на порядок
сильнее (медианное повторение \num{3.26} против \num{1.02} у контрольной
группы). Этот пилот — одно испытание на шести задачах, результатом его
считать нельзя. Честно про охват можно сказать так: наша находка держится
там, где окно не насыщено, а направление эффекта в насыщенном режиме
измерено, но не установлено.

\paragraph{Точность реализации.} Наш блок состояния на шаге толще, чем
сообщает оригинал\live, так что абсолютные величины стоимости — верхние
границы (\S\ref{sec:repro}). Обе поправки работают против нашего ключевого
результата.

\paragraph{Перенос тарифов в воспроизведении.} Раздел~\ref{sec:repro}
применяет наши измеренные цены к опубликованным числам токенов оригинала.
Порядок величины от точных множителей не зависит; они — оценка, и
помечены как таковая.

\paragraph{Контаминация, принципиально не обнаруживаемая.} Правило
карантина (\S\ref{sec:pitfalls}) опирается на ошибочные ходы. Деградацию
провайдера ниже порога видимости ошибок оно бы не поймало. Границу ущерба
задают два факта: результаты по стоимости вообще не зависят от вердиктов, а
результат по способности решать задачи пережил удаление каждой ячейки,
которую правило всё же поймало.

\paragraph{Две невыполнимые задачи.} Две из 89 задач отказывают при
настройке контейнера при каждом условии\live{} и отсутствуют во всех руках
одинаково. Это убирает две задачи из корпуса, но не даёт преимущества
какому-либо условию.

\paragraph{Исследование композиции несёт одно испытание.}
Раздел~\ref{sec:composition} приводит \VRuns{} запусков за одно испытание
против трёх для E1. С учётом собственной истории проекта (эффект, значимый
при двух испытаниях, умер на третьем (\S\ref{sec:pitfalls})), отдельные шаги
от ступени к ступени там нужно считать предварительными. Сквозное сравнение
и направление тренда стоимости согласуются на каждой ступени — поэтому они
вообще приводятся, — но ни одно утверждение в аннотации или заключении этой
статьи на них не опирается.

\paragraph{Нет руки только-хвоста.} Четыре условия E1 — это контрольная
группа без блока состояния и три руки, несущие состояние с разными
хвостами. Руки, которая несла бы последние $\tailk$ обменов и \emph{не}
несла бы блока состояния, среди них нет. Без неё \S\ref{sec:capability}
может показать, что восстановление истории устраняет дефицит, но не может
отделить, какая часть оставшегося поведения при $\tailk{=}6$ приходит от
блока состояния, а какая — от истории рядом с ним. Таблица~\ref{tab:totals}
подсказывает, что это важно: \StateK{6} завершает серию за меньшее общее
число ходов, чем контрольная группа, а чисто аддитивный учёт двух
компонентов такого исхода очевидным образом не предсказывает. Мы отмечаем
это как необъяснённое, не интерпретируя, и этой руке место в следующей
серии.
